\section{Discussion}

The method deliberately trades precision in individual analog components for system-level observability and digital assistance. Its principal benefit is not that the ADC output has zero mean; a software subtraction could accomplish that. Instead, the DACs modify the physical analog operating points before large downstream gains are applied. The correction is therefore largely agnostic to the source of the slowly varying error: amplifier offset, reference error, component variation, and other drift mechanisms are rejected by the same loop, provided that the affected stage is observable and the DAC has sufficient control authority and range.

The foreground method was adopted after implementing a continuously active DC servo. That experiment showed that lowering the servo corner enough to protect slow vibration components produced an excessively long centering transient. Increasing the servo bandwidth improved settling but interfered with the desired low-frequency signal. A programmable-gain change also modified the loop conditions and forced the system through another transient. In contrast, the proposed controller is disconnected and its DAC outputs are held constant throughout acquisition. It therefore adds neither time-varying feedback nor extra servo poles to the signal path, allowing the analog stages to be designed for the desired measurement band. EEPROM records also avoid the full calibration transient after most gain changes: stored codes are first restored and verified, and automatic recalibration is needed only when they no longer satisfy the tolerance. The remaining costs are a temporary interruption when recalibration is necessary, dependence on a sufficiently quiet input during that interval, and finite DAC resolution. The one-LSB refinement cannot remove a residual smaller than the change produced by one DAC code, and calibration cannot correct noise, distortion, loss of observability, or a stage that has exhausted the available correction range.

Only one board of the final four-VDAC revision is currently available. Repeated trials on this unit can establish algorithmic repeatability, but they do not characterize unit-to-unit process variation. The three earlier board revisions are useful for design history but are not statistically interchangeable with the final hardware.

